\subsection{Geometrische Definition}
\label{subsec:vector_algebra_scalar_product_definition}

\begin{DefinitionBox}{Skalarprodukt}{vector_algebra_scalar_product}
  Für zwei von null verschiedene Vektoren $\vec a,\vec b\in\R^n$ und den
  eingeschlossenen Winkel $0\le\varphi\le\pi$ ist das
  \emph{Skalarprodukt} \EnglishTerm{dot product} definiert durch
  \[
    \boxed{\vec a\cdot\vec b=|\vec a|\,|\vec b|\cos(\varphi)}.
  \]
  Für den Nullvektor wird $\vec a\cdot\vec 0=\vec 0\cdot\vec a=0$ gesetzt.
\end{DefinitionBox}

\noindent
\begin{minipage}[c]{0.49\linewidth}
  \centering
  \CachedTikzFigure{content/resources/chapters/vector_algebra/tikz/dot_product_geometry.tex}
\end{minipage}\hfill
\begin{minipage}[c]{0.47\linewidth}
  \begin{KeyPointBox}{Das Ergebnis ist ein Skalar}{vector_algebra_dot_product_scalar}
    \[
      \underbrace{\vec a}_{\text{Vektor}}
      \cdot
      \underbrace{\vec b}_{\text{Vektor}}
      =
      \underbrace{s}_{\text{Skalar}}\in\R.
    \]
    Das Skalarprodukt ist kein Vektor und wird nicht komponentenweise als
    Vektor ausgegeben.
  \end{KeyPointBox}
\end{minipage}

\begin{ExampleBox}{Skalarprodukt aus Längen und Winkel}{vector_algebra_geometric_dot_product_example}
  Für $|\vec a|=2$, $|\vec b|=3$ und $\varphi=30^\circ$ gilt
  \[
    \vec a\cdot\vec b
    =2\cdot3\cdot\cos(30^\circ)
    =\Cloze{4cm}{3\sqrt3}.
  \]
\end{ExampleBox}

\begin{KeyPointBox}{Vorzeichen und eingeschlossener Winkel}{vector_algebra_dot_product_sign}
  \begin{center}
    \CachedTikzFigure{content/resources/chapters/vector_algebra/tikz/dot_product_signs.tex}
  \end{center}
\end{KeyPointBox}
